拿到一个小工程时，先用熟悉的工具跑出一份可以对照的结果。第~\ref{chap:ide-tools}~章的计分器是最小的练习：用 Debug 构建停在函数调用，查看当前栈帧与变量，回答“程序在哪一步偏离预期”。

请 AI 帮忙时，把输入、预期和实际现象说清楚，再决定要它解释，还是让代理读取项目。按第~\ref{sec:ai-first-practice}~节的步骤先写下判断，再看修改、重新构建并用独立输入核对。工具操作随用随查；需要用 Git 审查、撤回或隔离改动时，从第~\ref{sec:ai-project-workflow}~节接着练。

当同一版本换到另一台机器却跑不通，接着问“运行条件哪里变了”。卷一第 4 章会带你在 Ubuntu 机载电脑上核对版本、依赖、工作目录和日志，把本地调试的办法接到实机部署上。
